\chapter{Chromosomes Bring Heredity into Cells}

\begin{kaichang}
Find a family portrait, or browse family photos on your phone. Notice the obvious first: you resemble both your father and your mother---perhaps his eye shape with her face shape. Now look at your siblings (or, if you have none, a photograph with your cousins): you come from the same parents, yet no two look exactly alike. One has double eyelids, another single; one has naturally curly hair, another straight. Even dimples may appear on only one face.

Now guess: why do the same parents, using the same ``recipe,'' make a different ``product'' each time? Is it an entirely random draw, or are there calculable rules? In Chapter 63, Mendel's peas showed that rules exist and can be written as ratios. Yet Mendel never saw what his ``hereditary factors'' looked like or where they hid in cells. This chapter gives those invisible, intangible factors a real home we can see under a microscope: chromosomes.
\end{kaichang}

\section{Where Do Mendel's Factors Live?}

Condense Chapter 63 into three statements: traits are determined by paired hereditary factors (now called alleles); when reproductive cells form, paired factors separate, each going its own way (segregation); factors for different traits usually combine independently (independent assortment). Growing plants and counting peas led Mendel to ratios as elegant as mathematical theorems.

But place yourself in a biologist's shoes around 1900. The theory had a gap large enough to keep you awake: \textbf{where are the factors}? Segregation says paired factors ``separate'' when reproductive cells form. That is a specific physical action: separation means they were together, then some force moved them to opposite sides. What exists in pairs inside a cell and separates precisely when reproductive cells form? Without such an object, Mendel's factors remain mathematical ghosts: the ratios work, but nobody knows what carries them out.

This chapter catches that ghost under the microscope. The clue comes from an apparently unrelated field: cell division.

\section{Zooming into the Nucleus}

During Chapter 57's cell tour, you saw the nucleus, the ``archive'' enclosed by a double membrane. It stores DNA, the carrier of all the cell's hereditary information. We will not rush into DNA's molecular structure (the subject of Chapters 65 and 66), but first consider its \textbf{packaging}.

Straighten all the DNA in an ordinary human cell and join it end to end, and it stretches about two meters---taller than you. Those two meters must fit inside a nucleus only about \ud{6}{\mu m} (six millionths of a meter) across. Imagine neatly packing a 200-kilometer thread into a table-tennis ball. Stuffing it in randomly would cause knots, and knotted thread cannot be read or copied. The cell winds DNA around protein spools called \textbf{histones}, then folds and bundles the wound spools in successive layers. This whole DNA-and-protein package normally spreads loosely, like a bowl of noodles, and is called \textbf{chromatin}. As division approaches, it condenses further into short, thick ``rods'' clearly visible under a microscope: the famous \textbf{chromosomes}. Remember: chromatin and chromosomes are not two different things, but the same material in different packaging states, like a quilt spread out normally and tied tightly for moving.

Each body cell (apart from a few exceptions such as mature red blood cells) contains 46 chromosomes in 23 pairs. One member of each pair came from your father's sperm, the other from your mother's egg. Such a pair consists of \textbf{homologous chromosomes}: they match in size and shape, with genes governing the same kinds of traits at corresponding positions---for example, an eye-pigment-related gene at a particular position. But ``the same kind of trait'' does not mean ``the same version'': the paternal chromosome might carry a ``double-eyelid version,'' the maternal one a ``single-eyelid version.'' Different versions at the same position are Chapter 63's alleles. Of the 23 pairs, 22 are the same in males and females and are called autosomes. The twenty-third pair consists of sex chromosomes: two X chromosomes in females, one X and a much smaller Y in males.

One last term is also the easiest to confuse. Before division, a cell makes a complete copy of each chromosome. The two identical copies do not immediately separate; like a zipper not yet opened, they remain joined at a central region called the \textbf{centromere}. These attached copies are \textbf{sister chromatids}. Keep the relationships clear: homologous chromosomes are a pair consisting of ``one from Dad, one from Mom''; sister chromatids are ``two photocopies of the same chromosome.'' Pairing and copying are entirely different. Soon you will see that meiosis derives its ingenuity from the order in which these two kinds of pair take the stage.

\section{Mitosis: Faithful Copying and Equal Shares}

Your body grows from one fertilized egg to about thirty-seven trillion cells by mitosis. Chapter 59 discussed cells cooperating to form a body; here, focus on division itself, whose logic is beautifully simple.

Mitosis solves this problem: when one cell becomes two, each new cell must receive a complete genetic archive. The solution has two steps. First, \textbf{copy before dividing}: duplicate all 46 chromosomes, each becoming two sister chromatids joined at the centromere. The cell now has 46 ``double-line'' chromosomes and twice the information. Second, \textbf{divide equally}: a protein-fiber ``traction system,'' the spindle, takes hold of chromosome centromeres from opposite ends, separates each joined sister pair, and pulls one copy left and one right. The cell then pinches into two.

The result? Forty-six chromosomes on the left and forty-six on the right, each from one member of a former sister pair. The daughter cells' archives match \textbf{letter for letter}. Skin cells producing skin cells, liver cells producing liver cells, and wound-edge cells dividing to fill a gap all use this mechanism: ``make a copy, distribute two equal shares.'' Mitosis maintains not a single cell but a lineage lasting decades. The archive in your liver cells today is one in a chain of copies of the archive present when you were three. (Copying occasionally goes wrong; that is Chapter 72's topic.)

Notice especially: throughout mitosis, homologous chromosomes mind their own business. The paternal and maternal versions never pair or exchange anything. Sister chromatids, not homologous chromosomes, are divided equally. The next section depends on this distinction.

\section{Meiosis: Halving, Shuffling, and Exchange}

If reproductive cells also arose by mitosis, trouble would follow: sperm and egg would each have 46 chromosomes; their child would have 92, its child 184, and so on. Doubling every generation would soon overflow the nucleus. Chromosome number needs a ``halving'' mechanism so that two ``half sets'' restore a complete set at fertilization. That mechanism is \textbf{meiosis}, which occurs only in cells producing sperm and eggs.

The rule is \textbf{copy once, divide twice}. The first division (meiosis I) separates homologous chromosomes---paternal and maternal versions enter different daughter cells. Only the second division (meiosis II) separates sister chromatids, much as mitosis does. After these successive divisions, each gamete (sperm or egg) contains just 23 chromosomes, one from each homologous pair. At fertilization, 23 from sperm meet 23 from egg, restoring exactly 46, stable generation after generation.

Now you know what carries out Mendel's law of segregation: ``paired factors separate'' means homologous chromosomes separate in meiosis I. Genes sit on chromosomes; when chromosomes separate, their genes follow. The ghost has a body.

Independent assortment is even more striking. At metaphase I, 23 homologous pairs line up in the cell's center. Each independently chooses ``which way Dad's chromosome faces and which way Mom's faces,'' like 23 independently tossed coins. Receiving the ``Dad version'' of chromosome 1 has no effect on which version of chromosome 2 you receive. This is independent assortment inside the cell: genes on different chromosomes really are shuffled independently.

\begin{qiaoliang}
Calculate this: considering only the random orientations of 23 chromosome pairs, how many different chromosome combinations can a person's gametes have? Each pair offers 2 choices, so 23 pairs give $2^{23}$, about 8.38 million. Your father can therefore produce roughly 8.38 million sperm types and your mother 8.38 million egg types. The chance you and a full sibling draw exactly the same combination is much lower than winning a lottery jackpot. This does not yet include the exchange discussed next, which raises the possible combinations far beyond any human population. The same parents can never produce two identical children (except identical twins, who arise by division of one fertilized egg). Everyone looking different in the family portrait is not an accident, but a mathematical necessity.
\end{qiaoliang}

Meiosis has one last and especially ingenious operation: \textbf{crossing over}. Early in meiosis I, paired homologous chromosomes press closely together. Then paternal and maternal versions break at corresponding positions, exchange a segment, and reconnect precisely. This is not accidental breakage: specialized enzymes carry out a precise operation. Each chromosome retains exactly the same number of genes, but now paternal and maternal segments share a chromosome. The chromosome entering a gamete is neither purely paternal nor purely maternal, but a new mixed version. Shuffling gains another level: not just whole chromosomes, but segments within them.

\begin{shiyan}
\textbf{Observe Dividing Cells (Microscope Required).} Materials: a school or commercially prepared permanent slide of ``onion root-tip mitosis'' and a light microscope. Steps: under low power, locate the root-tip growth region, where cells are small, densely packed, and roughly square. Switch to high power and scan slowly for cells with clearly visible chromosomes. Observations: some cells have chromosomes lined up along a central line, about to divide equally; others have two groups moving toward opposite ends; still others have intact nuclear membranes and loose chromatin, not currently dividing. Count: approximately what fraction of the cells are dividing? Safety: permanent slides are already fixed and stained, with no chemical risk, but glass edges may be sharp; hold them by the sides. Do not slice onion root tips with a kitchen knife and stain them at home. Common stains can be corrosive and toxic; prepare slides yourself only in a school laboratory under a teacher's supervision. This observation matters because you see firsthand a process occurring continually inside every multicellular organism for billions of years.
\end{shiyan}

\section{Putting Letters on Chromosomes}

Now the symbols can appear. In Chapter 63, you used A and a for alleles. We now know their address: each gene occupies a definite chromosome position, called its \textbf{locus}. Picture a chromosome as a long street and genes as street addresses. The paternal and maternal streets have corresponding addresses (homology), but different ``versions'' can live at each one.

In symbols, meiosis does this: a homologous pair, one labeled A and the other a, separates into different gametes. A gamete receives either A or a, each with probability one-half. For A/a and B/b on different chromosomes, the two pairs orient independently, giving AB, Ab, aB, and ab gametes in equal proportions. This matches Chapter 63's Punnett square perfectly, but now every symbol has a real ``rod'' behind it.

What if two genes occupy \textbf{the same} chromosome? Suppose A and B share a street, with the paternal street carrying ``AB'' and the maternal street ``ab.'' They no longer assort independently: without crossing over, the street passes on whole, giving only AB or ab gametes, not Ab or aB. This is \textbf{linkage}: two genes are ``tied'' to one rod and tend to be inherited together.

Crossing over partly unties them: an exchange between their loci produces new combinations such as Ab and aB. Farther-apart loci have more opportunities for exchange between them and yield more new combinations; closer loci are harder to separate. In other words, \textbf{how frequently two genes separate reveals their chromosome distance}. That modest-sounding statement is the entire principle behind humanity's first gene maps. The evidence story below shows how it happened.

\begin{wuqu}
``Crossing over is a cellular malfunction that breaks chromosomes.'' Quite the opposite: it is a \textbf{standard step} of meiosis. Each time sperm or eggs form, every homologous pair usually undergoes one to several exchanges. Specialized enzymes cut, align, and repair at precisely matching positions. A lack of crossing over can instead cause trouble: besides producing new combinations, exchanges help homologs hold securely ``hand in hand'' before separating. Too few exchanges are associated with increased nondisjunction risk (the aneuploidy discussed later). Calling crossing over a malfunction is like calling ``shuffling cards'' ``tearing cards.'' Shuffling keeps the game fair and fresh. Real accidents are different: broken chromosomes joined at the wrong position (translocation), or whole segments lost or duplicated. Those are structural variants, related to Chapter 72's mutations.
\end{wuqu}

\section{Evidence: Grasshoppers and a White-Eyed Fly}

\begin{zhengju}
In 1902--1903, two trails converged under the microscope. American Walter Sutton studied grasshopper reproductive cells, while German Theodor Boveri experimented with sea-urchin eggs. Independently, both noticed the same thing: chromosome behavior in meiosis precisely paralleled Mendel's factors. Factors pair, chromosomes pair; factors segregate, chromosomes segregate; factors assort independently, chromosomes orient independently; both restore pairs at fertilization. They proposed a bold hypothesis: \textbf{genes reside on chromosomes}. This became the ``chromosome theory of inheritance.'' But at this stage it was only a hypothesis: parallel behavior does not prove a shared location. A shadow moving with a person might still be coincidence.

Morgan and his fruit flies secured the hypothesis. In 1910, Morgan found a white-eyed male among his cultured red-eyed flies. He crossed it with a red-eyed female; all first-generation offspring had red eyes, matching Mendel with red dominant. Intercrossing that generation produced roughly 3:1 red to white in the second, again matching Mendel. But one detail would not fit: \textbf{almost all white-eyed flies were male}. If eye color's gene lay on an autosome, sex and eye color should have no connection. Morgan proposed the only self-consistent explanation: the gene lay on the X chromosome. Males have only one X, inherited from their mother, and express whatever version it carries. Females have two X chromosomes and need white-eye versions on both to be white-eyed, making white eyes much more common in males. For the first time, a gene was ``located'' on a specific chromosome. Hypothesis became tangible.

The story has a second half. Morgan's student Sturtevant noticed that genes on the same X chromosome varied in how tightly they traveled together. Some almost never separated; others often separated through crossing over. He recalled the modest principle above: separation frequency measures distance. In 1913, still an undergraduate, Sturtevant converted six genes' exchange frequencies into distances and drew history's first linkage map: a ``chromosome street map'' showing only relative positions and distances, yet correctly predicting the gene order later observed directly. Its greatness lay in using the ``plodding'' method of counting offspring to measure molecular distances invisible to any microscope of the day. Modern genome sequencing has, of course, recorded street addresses in detail, but the first map was drawn with paper, a pen, and fruit flies.
\end{zhengju}

\section{When the Machinery Fails: Aneuploidy}

Any mechanical process can go wrong, including meiosis. One major error is \textbf{nondisjunction}: a homologous pair (or sister chromatids) fails to separate when it should. One gamete then receives 24 chromosomes, with an extra copy of one chromosome, while another is one short. Fertilization involving the extra-copy gamete produces an embryo with 47 chromosomes: three copies of one chromosome instead of two. A chromosome number that is not an exact multiple of the set is called \textbf{aneuploidy}.

You have probably heard of one example: an extra chromosome 21, or trisomy 21, known as Down syndrome. It is among the commonest numerical chromosome abnormalities in newborns and brings varying developmental differences and health risks. Why is chromosome 21 most common? One possible reason is that other trisomies disrupt development more severely and most pregnancies cannot continue beyond very early stages. An extra copy of only the smallest chromosomes (21 is among the smallest autosomes) can still allow development to birth. This is not ``chromosome 21 makes more mistakes,'' but ``among mistakes compatible with birth, 21 predominates.'' Observed frequency has passed through two filters: error rate and survival.

Statistics also show that the probability of a child with trisomy 21 rises with maternal age. This relates to egg cells' unusual history: they begin meiosis I while the mother herself is a fetus, then \textbf{pause} for decades until ovulation allows completion to continue. Imagine the difficulty of maintaining chromosome ``handholds'' flawlessly for decades. Two points matter: this is a population-level probability pattern, not a statement about an individual; and fertilization involves both sperm and egg, with nondisjunction also possible during sperm formation. Consult a medical professional for relevant questions about yourself or relatives. A popular-science book does not provide individual diagnosis or treatment advice.

Aneuploidy's deeper lesson is that hereditary information requires not only ``correct content'' but ``correct copy number.'' Gene products resemble recipe ingredients: an extra or missing spoonful changes the balance of the whole system. Dosage is itself part of information.

\section{Back to the Family Portrait}

Return to the opening question: why can the same parents not produce two identical children?

You now have the whole chain. During sperm production, the father's 23 homologous pairs are randomly halved, giving $2^{23}$ combinations; the same happens during egg production. Fertilization joins two ``grand lotteries.'' Whole-chromosome assortment alone gives about $2^{23} \times 2^{23}$, or seventy trillion, combinations. Add crossing over within each chromosome---the chromosome 1 a father passes on is actually a new mixture of the paternal grandparents' two versions---and the number of combinations becomes practically nameless.

Everyone in the family portrait is therefore unique. Your paternal inheritance is a hand your father drew from his parents, then reshuffled through crossing over; your siblings receive another shuffle of the same deck. You share roughly half your alleles, hence the resemblance, but nobody draws exactly the same hand, hence the differences. Resemblance and difference are neither fate nor coincidence, but factory settings of meiosis's ``shuffling machine.''

This also answers a familiar question: ``Why do relatives keep arguing about which parent a child resembles more?'' At the chromosome level, the answer is unexciting: apart from the small exceptions of sex chromosomes and mitochondria, the autosomal shares from both parents are exactly equal. As for ``looking more like'' someone, that depends on which genes happen to be dominant, which traits stand out, and whether observers look first at noses or eyes. It is more a game of attention than genetic favoritism.

\section{What Carries Information in a Chromosome?}

We have housed Mendel's factors in chromosomes, apparently completing the story. But a reader who enjoys challenging assumptions should now ask: chromosomes package \textbf{two} things together---DNA and protein. Which is the actual hereditary material?

Do not rush to say DNA. From the standpoint of 1930, those betting on protein had good reasons. Proteins contain 20 amino-acid types, with enormous variety and countless folded shapes, seemingly well suited to storing complex information. DNA, by contrast, was thought to be a monotonous repetition of four nucleotides, an unremarkable structural material. Choosing the carrier of heredity cannot be settled by guessing: it takes experiments, especially elegant ones that separate the two molecular types cleanly and test each independently.

Who succeeded, and how? What arguments and misjudgments occurred, and what female scientist's contribution was undervalued? In Chapter 65, we turn to a famous scientific ``identity investigation'': who proved that DNA carries hereditary information?

\begin{tiaozhan}
Why could Sturtevant draw his gene map as a straight line? For nearby genes, exchange frequency is approximately proportional to distance: double the distance, roughly double the opportunity for exchange. This defines a ``map-distance'' unit: if an average of 1 in 100 offspring is a recombinant between two loci, their distance is 1 unit (now called 1 centimorgan, cM, honoring Morgan). Suppose A--B exchange frequency is 5\%, B--C is 3\%, and B lies between A and C. What is the approximate A--C frequency? Your first thought is $5\%+3\%=8\%$, but the observed value is slightly lower. An exchange in \textbf{each} interval A--B and B--C (a double crossover) restores the original combination at the A and C ends, escaping the count. The double-crossover probability is about $0.05 \times 0.03 = 0.15\%$, giving an observed value of roughly $8\% - 2\times0.15\% = 7.7\%$. This is how geneticists detected double crossing over and thereby confirmed genes' linear arrangement. At greater distances, proportionality becomes badly distorted and requires more sophisticated correcting functions. You can skip this without losing the thread, but whenever you see a ``gene map,'' you now know that every unit of distance was earned by counting thousands upon thousands of offspring.
\end{tiaozhan}

\section*{Try It Yourself}

\begin{sikaoti}
\item Draw a homologous pair: a paternal chromosome labeled A (dominant) and a maternal one labeled a. First draw their replicated forms with sister chromatids and labeled centromeres, then show what each of four gametes receives after meiosis I and II. Add one sentence: which step separates sister chromatids, and which separates homologous chromosomes?
\item Both mitosis and meiosis include ``sister chromatid separation,'' yet one ``maintains'' and the other ``halves.'' Starting at 46 chromosomes, state each process's initial and final counts. Why must meiosis first separate homologous chromosomes for sister separation to contribute to a halved set?
\item A couple's first child has the father's single eyelids, the second the mother's double eyelids. A relative says, ``The third is due to resemble Dad.'' Explain the error using random assortment in meiosis. If this trait were controlled by one allele pair and the mother were heterozygous, what would be the probability of single eyelids in the third child?
\item Genes A and B are linked on one chromosome; the father's version is AB and the mother's ab. Without crossing over, which combinations can their child receive from the father's sperm? Which additional combinations become possible with crossing over? If A and B are far apart and exchange frequently, what classic ratio do offspring combinations approach? What does this show about the relationship between linkage and independent assortment?
\item Draw a diagram of information flow, starting with ``a homologous chromosome pair,'' passing through meiosis and fertilization, and ending with ``the child's 46 chromosomes.'' Mark where information is halved, where it is recombined, and where the double set is restored.
\item Cells of a person with trisomy 21 contain 47 chromosomes. Describe two possible nondisjunction events producing this outcome, one in meiosis I and one in meiosis II. Why cannot counting chromosomes under a microscope distinguish which division contained the error?
\end{sikaoti}
